% ============================================================
%  RESUME / ABSTRACT
% ============================================================
\chapter*{Abstract}
\addcontentsline{toc}{chapter}{Abstract}

The widespread adoption of mobile technologies has led to an exponential growth of
health applications worldwide. According to the World Health Organization, more than
7~billion mobile phone subscriptions exist globally, and in many countries access to
wireless technologies is easier to obtain than access to clean water or electricity.
In healthcare systems such as the Tunisian system, however, patient medical records
remain largely fragmented and undigitised. Patients carry disorganised paper folders
across multiple practitioners and facilities, leading to redundant tests, overlooked drug
interactions, and in the most serious cases, overdoses. This situation is especially
critical for pregnant women, where every medical decision affects both the mother and
the unborn child.

Sahhty, meaning ``my health'' in Tunisian Arabic, is a mobile health application built to
address these challenges. It allows patients to record and track their health measurements
and to follow their pregnancy week by week. After each new measurement, an artificial
intelligence engine evaluates the patient's risk level and sends an immediate alert when
intervention is needed. The platform integrates a national medication database built from
the DPM medication catalogue, the CNAM reimbursement list, the Thériaque drug
interaction database, and the LECRAT teratogenic risk classifications, allowing patients
to verify whether a treatment is safe during pregnancy and to detect incompatibilities
with their known allergies.

The application also allows patients to find nearby doctors on an interactive map, book
appointments, and store personal medical documents with controlled sharing. Doctors
have a dedicated space to manage their schedule and consult authorized patient files.
Real-time notifications keep both parties informed at all times. The platform supports
French and English and adapts its interface to the user's profile.

\bigskip
\noindent\textbf{Keywords:} Maternal health, Pregnancy monitoring, Mobile application,
Artificial intelligence, Drug interactions, Medications, Medical file, Tunisia.

\newpage
% ── French Abstract ────────────────────────────────────────────
\chapter*{Résumé}
\addcontentsline{toc}{chapter}{Résumé}

L'adoption généralisée des technologies mobiles a entraîné une croissance exponentielle
des applications de santé dans le monde. Selon l'Organisation Mondiale de la Santé,
plus de 7~milliards d'abonnements mobiles existent à l'échelle mondiale, et dans de
nombreux pays, l'accès aux technologies sans fil est plus facile que l'accès à l'eau
potable ou à l'électricité. Dans des systèmes de santé comme le système tunisien,
cependant, les dossiers médicaux des patients restent largement fragmentés et non
numérisés. Les patients transportent des dossiers papier désorganisés d'un praticien à
l'autre, ce qui entraîne des examens redondants, des interactions médicamenteuses
ignorées, et dans les cas les plus graves, des surdosages. Cette situation est
particulièrement critique pour les femmes enceintes, car chaque décision médicale
engage à la fois la santé de la mère et celle de l'enfant à naître.

Sahhty, qui signifie ``ma santé'' en arabe tunisien, est une application mobile de santé
conçue pour répondre à ces défis. Elle permet à la patiente d'enregistrer ses mesures de
santé et de suivre sa grossesse semaine par semaine. Après chaque nouvelle mesure, un
moteur d'intelligence artificielle évalue le niveau de risque et envoie immédiatement
une alerte si une intervention est nécessaire. La plateforme intègre une base de
médicaments nationale construite à partir du catalogue de la DPM, de la liste de
remboursement de la CNAM, de la base d'interactions médicamenteuses Thériaque, et
des classifications du risque tératogène LECRAT, permettant à la patiente de vérifier si
un traitement est sans danger pendant la grossesse et de détecter les incompatibilités
avec ses allergies connues.

L'application permet également de trouver des médecins à proximité sur une carte
interactive, de prendre des rendez-vous en ligne, et de stocker des documents médicaux
avec un contrôle total sur les accès. Les médecins disposent d'un espace dédié pour
gérer leur planning et consulter les dossiers des patientes autorisées. Des notifications
en temps réel tiennent les deux parties informées à tout moment. La plateforme est
disponible en français et en anglais et adapte son interface au profil de l'utilisateur.

\bigskip
\noindent\textbf{Mots-clés~:} Santé maternelle, Suivi de grossesse, Application mobile,
Intelligence artificielle, Interactions médicamenteuses, Médicaments, Dossier médical,
Tunisie.
